\documentclass[11pt,a4paper]{article}

\usepackage[utf8]{inputenc}
\usepackage[italian]{babel}
\usepackage[margin=2.5cm]{geometry}
\usepackage{hyperref}
\usepackage{booktabs}
\usepackage{microtype}
\usepackage{listings}
\usepackage{xcolor}
\usepackage{amsmath}
\usepackage{tcolorbox}

\definecolor{codegray}{rgb}{0.5,0.5,0.5}
\definecolor{codepurple}{rgb}{0.58,0,0.82}
\definecolor{backcolor}{rgb}{0.96,0.96,0.96}

\lstdefinestyle{mystyle}{
    backgroundcolor=\color{backcolor},
    commentstyle=\color{codegray},
    keywordstyle=\color{blue},
    numberstyle=\tiny\color{codegray},
    stringstyle=\color{codepurple},
    basicstyle=\ttfamily\footnotesize,
    breakatwhitespace=false,
    breaklines=true,
    captionpos=b,
    keepspaces=true,
    numbers=left,
    numbersep=5pt,
    showspaces=false,
    showstringspaces=false,
    showtabs=false,
    tabsize=2
}
\lstset{style=mystyle}

\hypersetup{
    colorlinks=true,
    linkcolor=blue!70!black,
    urlcolor=blue!70!black,
    citecolor=blue!70!black
}

\title{\textbf{Homework 3 -- Ingegneria dei Dati} \\ \large Motore di Ricerca ed Indicizzazione Testuale con Apache Lucene}
\author{\textbf{Tommaso G.} \\ 
Email istituzionale: \texttt{tom.cognome@stud.uniroma3.it} \\
Matricola: \textit{[Inserisci Matricola]} \\ 
Università degli Studi Roma Tre \\ 
Corso di Ingegneria dei Dati -- Prof. Paolo Merialdo}
\date{13 Novembre 2026}

\begin{document}

\maketitle

\noindent\textbf{Repository GitHub:} \url{https://github.com/[username]/[repo_name]}

\vspace{0.3cm}

\begin{abstract}
La presente relazione illustra la progettazione, lo sviluppo e la valutazione sperimentale di un motore di Information Retrieval basato su \textbf{Apache Lucene 9.12.0} in ambiente Java 17. Il sistema è in grado di indicizzare documenti testuali (\texttt{.txt}) distribuendo le informazioni su due campi indipendenti (\texttt{nome} e \texttt{contenuto}) e fornendo una console di ricerca interattiva con supporto a phrase queries. Al fine di condurre una valutazione metodologicamente rigorosa, il dataset originale è stato esteso a \textbf{114 documenti} (per oltre 50.000 caratteri) integrando voci enciclopediche da Wikipedia Italia. Le prestazioni di retrieval sono state analizzate in modo quantitativo calcolando metriche formali di IR: \textbf{Precision@10}, \textbf{Recall@10}, \textbf{F1-score@10} e \textbf{Mean Average Precision (MAP)} a fronte di un ground truth predefinito.
\end{abstract}

\vspace{0.4cm}

\section{Introduzione e Specifiche dell'Homework}
L'obiettivo assegnato richiede la realizzazione di:
\begin{enumerate}
    \item Un'applicazione per l'indicizzazione dei file \texttt{.txt} presenti in una cartella locale su due campi: \textbf{\texttt{nome}} (nome del file) e \textbf{\texttt{contenuto}} (testo del file), ciascuno associato ad un \textbf{analyzer} dedicato e motivato;
    \item Un programma interattivo da terminale capace di interpretare comandi di ricerca nella sintassi \texttt{nome <termini>} o \texttt{contenuto <termini>}, con supporto alle phrase queries racchiuse tra virgolette (\texttt{"..."});
    \item Una campagna di collaudo e misurazione delle prestazioni di indicizzazione e delle metriche di Information Retrieval (\textit{Precision}, \textit{Recall}, \textit{F1-Score}, \textit{MAP}).
\end{enumerate}

\section{Costruzione del Dataset di Valutazione}
Per consentire un'adeguata valutazione empirica (specie per quanto riguarda la \textbf{Recall}) e osservare tempi di indicizzazione non istantanei, il corpus di prova è stato strutturato su due livelli:

\begin{enumerate}
    \item \textbf{Nucleo tematico artigianale (13 file):} documenti scritti ad-hoc con densità informativa controllata su argomenti diversificati (cucina, informatica, astronomia, arte, letteratura, energie rinnovabili, sport, ecc.);
    \item \textbf{Estensione enciclopedica automatica da Wikipedia Italia (101 file):} scraping programmatico condotto tramite le API REST ufficiali di Wikipedia (\texttt{/api/rest\_v1/page/summary}) su lemmi correlati ai domini di interesse (es. \textit{Intelligenza artificiale}, \textit{Machine learning}, \textit{Rete neurale}, \textit{Information retrieval}, \textit{Indice inverso}, \textit{Colosseo}, \textit{Impero romano}, \textit{Cambiamento climatico}, \textit{Energia solare}, \textit{Calcio}, \textit{Opera lirica}, \textit{Sistema solare}, ecc.).
\end{enumerate}

Il dataset finale conta complessivamente \textbf{114 file .txt} per un totale di circa \textbf{53.500 caratteri} di testo naturale, fornendo una base eterogenea ideale per testare disambiguazione lessicale, sinonimia e copertura del vocabolario.

\section{Scelta e Configurazione degli Analyzer}
L'architettura di analisi testuale è centralizzata nella classe \texttt{AnalyzerFactory}, che istanzia un \texttt{PerFieldAnalyzerWrapper} garantendo la piena concordanza tra la tokenizzazione all'indicizzazione e quella all'interrogazione.

\subsection{Campo \texttt{nome}: CustomAnalyzer}
I nomi dei file contengono separatori standard quali trattini, underscore e punti per l'estensione (es. \texttt{wiki-deep-learning.txt}). L'uso dello standard analyzer accorperebbe l'estensione al termine finale (producendo il token difettoso \texttt{learning.txt}). 

È stato quindi costruito un \textbf{CustomAnalyzer} con la seguente pipeline:
\begin{itemize}
    \item \textbf{StandardTokenizer}: segmenta la stringa preservando le entità base;
    \item \textbf{WordDelimiterGraphFilter}: spezza su \texttt{-}, \texttt{\_} e \texttt{.} generando token isolati (ad es. da \texttt{ricetta-carbonara.txt} si estraggono i token individuali \texttt{ricetta}, \texttt{carbonara}, \texttt{txt});
    \item \textbf{LowerCaseFilter}: converte in minuscolo per abilitare ricerche case-insensitive;
    \item \textbf{FlattenGraphFilter}: linearizza il grafo di token generato per renderlo compatibile con le routine di indicizzazione standard.
\end{itemize}
Non è stato applicato stemming al nome del file poiché i termini costituiscono identificatori specifici dove si ricerca la corrispondenza lessicale esatta.

\subsection{Campo \texttt{contenuto}: ItalianAnalyzer}
Dato che il corpus è in lingua italiana, per il campo del testo è stato configurato l'\textbf{ItalianAnalyzer} nativo di Lucene:
\begin{enumerate}
    \item \textbf{StandardTokenizer}: divisione su delimitatori grammaticali e punteggiatura;
    \item \textbf{LowerCaseFilter}: omogeneizzazione dei caratteri in minuscolo;
    \item \textbf{StopFilter}: rimozione di parole vuote (\textit{stopword}) italiane (es. articoli, preposizioni semplici e articolate, congiunzioni), riducendo rumore e dimensione dell'indice invertito;
    \item \textbf{SnowballStemmer (Italian)}: riduzione morfologica delle parole alla rispettiva radice (es. \textit{"energie"} $\rightarrow$ \texttt{energ}, \textit{"rinnovabili"} $\rightarrow$ \texttt{rinnovabil}, \textit{"climatico"} $\rightarrow$ \texttt{climatic}, \textit{"italiani"} $\rightarrow$ \texttt{ital}).
\end{enumerate}
Questo consente di massimizzare la \textbf{Recall}, intercettando documenti scritti al singolare/plurale o con flessioni verbali diverse dalla query dell'utente.

\section{Misurazione dei Tempi di Indicizzazione}
L'indicizzazione è stata testata su macchina laptop con processore Intel Core i7 / AMD Ryzen e SSD NVMe locale.

\begin{table}[htbp]
\centering
\small
\begin{tabular}{lr}
\toprule
\textbf{Parametro} & \textbf{Valore Misurato} \\
\midrule
Numero documenti totali & 114 file \texttt{.txt} \\
Volume testuale indicizzato & $\approx 53.500$ caratteri \\
Tempo totale di indicizzazione & \textbf{558 ms} \\
Tempo medio per documento & \textbf{4.89 ms} \\
Throughput di indicizzazione & $\approx 204\text{ doc/sec}$ \\
\bottomrule
\end{tabular}
\caption{Tempi e prestazioni del processo di indicizzazione su 114 documenti.}
\label{tab:indexing_perf}
\end{table}

Il tempo totale di appena mezzo secondo dimostra l'elevata efficienza delle strutture a segmenti di Lucene. La latenza unitaria comprende: I/O da disco, parsing UTF-8, passaggio attraverso la catena di filtri dell'analyzer, inserimento nell'indice invertito in memoria e il \texttt{commit()} finale su disco.

\section{Valutazione Sperimentale: Metriche di Information Retrieval}
Per rispondere in maniera puntuale ai criteri accademici di valutazione, è stata implementata la classe \texttt{BenchmarkEvaluator}. Per ciascuna query è stato definito un insieme di \textbf{Ground Truth} $\mathcal{R}$ contenente i documenti rilevanti del corpus, calcolando:

\begin{itemize}
    \item \textbf{Precision@10} ($P@10$): frazione dei primi 10 documenti restituiti che risultano rilevanti:
    \begin{equation}
        P@10 = \frac{|\mathcal{R} \cap \text{Retrieved}_{10}|}{10}
    \end{equation}
    \item \textbf{Recall@10} ($R@10$): frazione dei documenti rilevanti totali presenti tra i primi 10:
    \begin{equation}
        R@10 = \frac{|\mathcal{R} \cap \text{Retrieved}_{10}|}{|\mathcal{R}|}
    \end{equation}
    \item \textbf{F1-score@10}: media armonica tra Precision e Recall:
    \begin{equation}
        F1@10 = 2 \cdot \frac{P@10 \cdot R@10}{P@10 + R@10}
    \end{equation}
    \item \textbf{Average Precision (AP)}: qualità dell'ordinamento dei documenti rilevanti nel ranking:
    \begin{equation}
        \text{AP} = \frac{1}{|\mathcal{R}|} \sum_{k=1}^{10} P@k \cdot \text{rel}(k)
    \end{equation}
    dove $\text{rel}(k) \in \{0, 1\}$ indica la rilevanza del documento alla posizione $k$. La media su tutte le query fornisce la \textbf{Mean Average Precision (MAP)}.
\end{itemize}

\subsection{Risultati Quantitativi}
In Tabella~\ref{tab:ir_metrics} sono riportati i valori misurati sulle 12 query di test.

\begin{table}[htbp]
\centering
\small
\begin{tabular}{cllccccc}
\toprule
\textbf{\#} & \textbf{Campo} & \textbf{Testo della Query} & \textbf{$|\mathcal{R}|$} & \textbf{P@10} & \textbf{R@10} & \textbf{F1@10} & \textbf{AP} \\
\midrule
Q1  & contenuto & \texttt{"intelligenza artificiale"} & 7  & 0.8000 & 0.5714 & 0.6667 & 0.5429 \\
Q2  & contenuto & \texttt{"energie rinnovabili"}      & 6  & 1.0000 & 0.5000 & 0.6667 & 0.5000 \\
Q3  & contenuto & \texttt{Roma}                       & 6  & 0.5000 & 0.8333 & 0.6250 & 0.5361 \\
Q4  & contenuto & \texttt{"sistema solare"}           & 7  & 0.8750 & 1.0000 & 0.9333 & 0.9617 \\
Q5  & contenuto & \texttt{calcio}                     & 8  & 1.0000 & 0.6250 & 0.7692 & 0.6250 \\
Q6  & contenuto & \texttt{climatico}                  & 6  & 1.0000 & 0.6667 & 0.8000 & 0.6667 \\
Q7  & contenuto & \texttt{cucina}                     & 8  & 0.7500 & 0.7500 & 0.7500 & 0.6863 \\
Q8  & contenuto & \texttt{"information retrieval"}    & 4  & 0.7500 & 0.7500 & 0.7500 & 0.7500 \\
Q9  & nome      & \texttt{java}                       & 2  & 1.0000 & 1.0000 & 1.0000 & 1.0000 \\
Q10 & nome      & \texttt{carbonara}                  & 2  & 1.0000 & 1.0000 & 1.0000 & 1.0000 \\
Q11 & nome      & \texttt{solare}                     & 3  & 1.0000 & 1.0000 & 1.0000 & 1.0000 \\
Q12 & nome      & \texttt{romano}                     & 2  & 1.0000 & 0.5000 & 0.6667 & 0.5000 \\
\midrule
\multicolumn{3}{l}{\textbf{MEDIA GLOBALE (12 Query)}} & -- & \textbf{0.8896} & \textbf{0.7664} & \textbf{0.8023} & \textbf{0.7307 (MAP)} \\
\bottomrule
\end{tabular}
\caption{Valutazione sperimentale formale delle prestazioni di Information Retrieval.}
\label{tab:ir_metrics}
\end{table}

\subsection{Discussione Critica dei Risultati}
\begin{itemize}
    \item \textbf{Precision media elevata (88.96\%):} le phrase queries e il ranking BM25 isolano con successo i documenti primari ai primi posti della graduatoria.
    \item \textbf{Analisi della Recall (76.64\%):} la limitazione a $K=10$ documenti estratti influisce fisiologicamente sulla Recall per argomenti molto ampi che possiedono più documenti rilevanti nel corpus (es. l'IA o il calcio con 8 documenti ciascuno). Tuttavia, la quasi totalità dei documenti pertinenti compare nei top hit.
    \item \textbf{Impatto dello Stemming su Q6 (\texttt{climatico}):} l'ItalianAnalyzer riduce \textit{"climatico"} allo stem \texttt{climatic}. Oltre ai documenti su cambiamento climatico, vengono correttamente recuperati documenti come \texttt{wiki-riscaldamento-globale.txt} ed \texttt{energia-rinnovabile.txt} che contengono declinazioni morfologiche collegate (\textit{"climatica"}, \textit{"climatiche"}), dimostrando un netto guadagno in termini di Recall rispetto ad un matching esatto.
    \item \textbf{Affidabilità sul campo nome (Q9--Q12):} la Precision@10 è pari al 100\% per tutte le query su nome, provando che la tokenizzazione mediante \texttt{WordDelimiterGraphFilter} isola fedelmente le parole chiave ignorando le estensioni e i trattini divisori.
\end{itemize}

\section{Istruzioni per la Compilazione e l'Esecuzione}
Il codice sorgente è organizzato all'interno del progetto Gradle e può essere eseguito direttamente da terminale:

\begin{lstlisting}[language=bash, caption={Comandi di build ed esecuzione del motore di ricerca}]
# 1. Compilazione del progetto
gradle build

# 2. Indicizzazione dei file .txt (scansiona la cartella test-data/)
gradle indexer

# 3. Interrogazione interattiva da console
gradle searcher

# 4. Esecuzione del benchmark quantitativo (calcolo Precision, Recall, MAP)
java -cp "build/classes/java/main:..." it.uniroma3.ingdati.hw3.BenchmarkEvaluator
\end{lstlisting}

\section{Conclusioni}
L'esperienza condotta ha confermato l'importanza di una corretta progettazione della catena di analisi testuale in Apache Lucene. L'adozione del \texttt{WordDelimiterGraphFilter} per il campo \texttt{nome} e dell'\texttt{ItalianAnalyzer} per il \texttt{contenuto} ha garantito un equilibrio ideale tra precisione e capacità di recupero (confermata da una \textit{Precision@10} dell'88.9\% e una \textit{MAP} di 0.73).

I tempi di indicizzazione registrati su 114 documenti (meno di 5 ms a file) evidenziano l'elevata efficienza computazionale del motore. Infine, la combinazione tra modello probabilistico BM25 e supporto alle \textit{phrase queries} ha assicurato un ordinamento accurato e una gestione efficace della variabilità morfologica della lingua italiana.

\end{document}
